\documentclass[12pt]{article}
\usepackage{amsmath}
\usepackage{amssymb}
\usepackage{amsgen}
\usepackage{amsfonts}

\newtheorem{theorem}{Theorem} 
\newtheorem{corollary}[theorem]{Corollary} 
\newtheorem{conjecture}[theorem]{Conjecture} 
\newtheorem{lemma}[theorem]{Lemma} 
\newtheorem{proposition}[theorem]{Proposition} 
\newtheorem{definition}[theorem]{Definition} 
\newtheorem{example}[theorem]{Example} 
\newtheorem{exercise}[theorem]{Exercise} 
\newtheorem{axiom}{Axiom} 
\newtheorem{remark}{Remark} 

\numberwithin{equation}{section}

\def\ppp{{\mathbb{P}}}
\def\aaa{{\mathbb{A}}}
\def\fff{{\mathbb{F}}}
\def\qqq{\mathbb{Q}}
\def\rrr{\mathbb{R}}
\def\ccc{\mathbb{C}}
\def\zzz{\mathbb{Z}}
\def\nnn{\mathbb{N}}
\def\pf{\noindent {\bf Proof}:\ }
\def\qed{$\Box$}

\begin{document}

\author{David Joyner\thanks{Mathematics Dept, 
US Naval Academy, Annapolis, MD 21402,
wdj@usna.edu}}

\title{Conjectural permutation decoding of some AG codes\thanks{MSC 2000: 
14H37, 94B27,20C20,11T71,14G50,05E20,14Q05}}

\maketitle

\begin{abstract}
We study the action of a finite group on the Riemann-Roch
space of certain divisors on a specific hyperelliptic curve
$X$ defined over a finite field with ``large'' automorphism group
$G$. If $D$ is a divisor on $X$ left stable by $G$ then 
the natural representation of $G$ on the Riemann-Roch space 
$L(D)=L_X(D)$. If $E=P_1+...+P_n$ is another
divisor on $X$ left stable by $G$ then this representation on
$L(D)$ induces a permutation representation on the
associated AG code $C=C(D,E)$ (defined below).
This note discusses the permutation decoding of these AG codes.
The main ``results'' are conjectures regarding the complexity of 
the permutation decoding of these hyperelliptic codes. 
\end{abstract}


\section{Introduction}

Here is the rough idea discussed in this note.

Let $X$ be a plane curve defined over a finite field
$F$ by means of some polynomial equation $f(x,y)=0$.
(Though strictly speaking we are identifying,
by abuse of notation, $X$ with the associated projective 
algebraic curve.) 
A {\it divisor} $D$ on $X$ is a formal integer linear
combination of elements in the set of $F$-rational
points $X(F)$, $D=d_1Q_1+...d_rQ_r$. The Riemann-Roch
space $L(D)$ is a vector space of rational functions 
in $x,y$ on $X$ whose zero/pole at $Q_i$ is
constroled by $d_i$, $1\leq i\leq r$. (See Stichtenoth
\cite{Sti} or Hartshorne \cite{H} for a precise definition.)
Let $E=P_1+...+P_n$ be a sum of points distinct from the
$Q_i$'s. The {\it AG code} $C=C_X(D,E)$ associated to $X$ is 
defined by 

\begin{equation}
\label{eqn:AGcode}
C=\{(f(P_1),...,f(P_n))\ |\ f\in L(D)\}.
\end{equation}
This is the image of $L(D)$ under the evaluation map

\begin{equation}
\label{eqn:eval}
\begin{array}{c}
eval_E:L(D)\rightarrow F^n,\\
f \longmapsto (f(P_1),...,f(P_n)).
\end{array}
\end{equation}
These codes were originally introduced by Goppa in the 
1980s and have proven to be a rich source of 
efficient error-correcting block codes.

This paper investigates one small class of such AG codes
which seems to have an unusually large amount of
symmetry. The hope is that this symmetry can be used to 
create a class of codes which has a very efficient
decoding procedure.

Some of these ideas are part of joint research with Amy Ksir and Will Traves.
For related research, see Joyner and Ksir \cite{JK1}.
See \cite{JK2} and \cite{JT} for further background.

\section{AG codes associated to $y^2=x^p-x$}

Properties of $y^2=x^p-x$ are discussed for example in 
G\"ob \cite{G}.

\subsection{Background on the curve}

In general, if $X$ is a curve defined over a field $F$
with finite automorphism group $G=Aut_F(X)$ then we call
$G$ {\bf large} if $|G|>|X(F)|$.

\begin{lemma}
\label{lemma:large}
If $G$ is large then every point of $X(F)$
is ramified for the covering $X\rightarrow X/G$.
\end{lemma}

\pf
Suppose $P\in X(F)$ is not ramified, so the stabilizer of $P$,
$G_P$, is trivial. In this case,
$|G\cdot P|= |G|/|G_P|= |G|$. But $G\cdot P\subset X(F)$
so $|G\cdot P|\leq |X(F)|$, a contradiction. \qed

\subsubsection{Case $F=GF(p)$}
\label{sec:F=GF(p)}

Let $p\geq 3$ be a prime,
$F=GF(p)$, and let $X$ denote the curve defined by

\[
y^2=x^p-x.
\]
This has genus $\frac{p-1}{2}$. 
We assume that the automorphism group $G=Aut_F(X)$ is
a central 2-fold cover of $PSL(2,p)$,
we have a short exact 
sequence,

\begin{equation}
\label{eqn:*}
1\rightarrow Z \rightarrow G\rightarrow PSL_2(p)\rightarrow 1,
\end{equation}
where $Z$ denotes the center of $G$ 
($Z$ is generated by the hyperelliptic involution).
The following transformations are elements of
$G$:

\begin{equation}
\label{eqn:3.2}
\begin{array}{cc}
\gamma_1=
\left\{
\begin{array}{c}
x\longmapsto x,\\
y\longmapsto -y,
\end{array}
\right. ,
&
\gamma_2=\gamma_2(a)=
\left\{
\begin{array}{c}
x\longmapsto a^2x,\\
y\longmapsto ay,
\end{array}
\right.\\
\gamma_3=
\left\{
\begin{array}{c}
x\longmapsto x+1,\\
y\longmapsto y,
\end{array}
\right. ,
&
\gamma_4=
\left\{
\begin{array}{c}
x\longmapsto -1/x,\\
y\longmapsto y/x^{\frac{p+1}{2}},
\end{array}
\right. 
\end{array}
\end{equation}
where $a\in F^\times$ is a primitive $(p-1)-st$ root of unity.
This group acts transitively on $X(F)$,
so it has an orbit of size $d_0=|X(F)|=p+1$.

Let $P_1=(1:0:1)$ and let $H$ be its stabilizer in $G$.
A counting argument shows that $H$ is a solvable group of order 
$2p(p-1)$ generated by 
$\gamma_1$, $\gamma_2(a)$ and $\gamma_3$.
By Lemma \ref{lemma:large}, every point in 

\[
X(F)=\{(1:0:0),(0:0:1),(1:0:1),...,(p-1:0:1)  \}
\]
is ramified over the covering $X\rightarrow X/G$
in the sense that each stabilizer $G_P=Stab_G(P)$
is non-trivial, $P\in X(F)$. 

It is known (Proposition VI.4.1, \cite{Sti}) that, 
for each $m\geq 1$, the Riemann-Roch space
of $D=mP_1$ has a basis consisting of monomials,

\[
x^iy^j,\ \ \ \ \ \ 
0\leq i\leq p-1,\ j\geq 0,\ 2i+pj\leq m.
\]

\begin{lemma}
The semisimplification $\rho_{ss}$ of 
the representation $\rho$ of $H$ acting on $L(D)$ 
is the direct sum of 
one-dimensional representations of $G$.
\end{lemma}

\pf The generator $\gamma_1$ acts trivially on the 
basis of $L(D)$, whereas
\[
\gamma_2(a):
\left(
\begin{array}{c}
1\\
x\\
\vdots \\
x^ry^s
\end{array}
\right)
\longmapsto 
\left(
\begin{array}{c}
1\\
a^2x\\
\vdots \\
a^{2r+s}x^ry^s
\end{array}
\right)
=
\left(
\begin{array}{cccc}
1 & 0 & ...& 0\\
0 & a^2 & ...& \vdots\\
\vdots &\ddots  & &0 \\
0 & ... & 0 & a^{2r+s}
\end{array}
\right)
\left(
\begin{array}{c}
1\\
x\\
\vdots \\
x^ry^s
\end{array}
\right),
\]
and

\[
\gamma_3:
\left(
\begin{array}{c}
1\\
x\\
\vdots \\
x^ry^s
\end{array}
\right)
\longmapsto 
\left(
\begin{array}{c}
1\\
x+1\\
\vdots \\
(x+1)^ry^s
\end{array}
\right)
=
\left(
\begin{array}{cccc}
1 & 0 & ...& 0\\
1 & 1 & ...& \vdots\\
\vdots &\ddots  & &0 \\
0 & ... & r & 1
\end{array}
\right)
\left(
\begin{array}{c}
1\\
x\\
\vdots \\
x^ry^s
\end{array}
\right),
\]
where the non-zero terms in bottom row of the matrix representation of
$\gamma_3$ are in the last $r+1$ row entries and
consist of the binomial coefficients $\frac{r!}{(r-j)!j!}$,
$0\leq j\leq r$.
Therefore, the group generated by these matrices is 
lower-triangular, hence solvable. \qed

\subsubsection{Case $F=GF(p^2)$}
\label{sec:gfp2}

Let $F=GF(p^2)$ and let $F_0=GF(p)$.

The automorphism group $G=Aut_F(X)$ is
a central 2-fold cover of $PGL(2,p)$ andwe have a short exact sequence,

\begin{equation}
\label{eqn:**}
1\rightarrow Z \rightarrow G
\stackrel{\tau}{\rightarrow} PGL_2(p)\rightarrow 1,
\end{equation}
where $Z$ denotes the subgroup of $G$ generated by the hyperelliptic
involution (which coincides with the center of $G$), 
by G\"ob \cite{G}.
The group $G$ has order $2|PGL(2,p)|=2p(p^2-1)$.
The following transformations generate $G$:

\[
\begin{array}{cc}
\gamma_1=
\left\{
\begin{array}{c}
x\longmapsto x,\\
y\longmapsto -y,
\end{array}
\right. ,
&
\gamma_2=\gamma_2(a)=
\left\{
\begin{array}{c}
x\longmapsto a^2x,\\
y\longmapsto ay,
\end{array}
\right.\\
\gamma_3=
\left\{
\begin{array}{c}
x\longmapsto x+1,\\
y\longmapsto y,
\end{array}
\right. ,
&
\gamma_4=
\left\{
\begin{array}{c}
x\longmapsto -1/x,\\
y\longmapsto y/x^{\frac{p+1}{2}},
\end{array}
\right. 
\end{array}
\]
where $a\in F^\times$ is a primitive $2(p-1)-st$ root
of unity.


\begin{proposition}
\label{prop:orbits} 
Let $p>3$ be a prime.
\begin{itemize}
\item[(a)] Case $p\equiv 3\, ({\rm mod}\, 4)$:

Let $P_1=(1:0:1)$ and fix some $P_2\in X(F)-X(F_0)$.
The set of rational points $X(F)$ decomposes into
a disjoint union

\[
C_1=X(F_0)=G\cdot P_1,\ \ \ \ \ \ 
C_2=X(F)-X(F_0)=G\cdot P_2,
\]
with $|C_1|=p+1$ and $|C_2|=2p(p-1)$.

\item[(b)] Case $p\equiv 1\, ({\rm mod}\, 4)$:

The automorphism group of $X/F$ acts transitively
on $X(F)$ and the stabilizer of any point is
a group of order $2p(p-1)$.

\end{itemize}
\end{proposition}

\begin{remark}
The proof of this proposition
is omitted, so may be regarded as a conjecture
instead, if the reader wishes.
It has been verified using MAGMA if $p=5, 7,11, 13$. 
It has been proven in an
email to the author by Bob Guralnick.
\end{remark}


This and Lemma \ref{lemma:large} imply 
every point in $X(F)$ is ramified for the covering
$X\rightarrow X/G$.

Let $P_1,P_2$ be as above and 
let $H_i$ be the stabilizer of $P_i$ in $G$ (for $i=1,2$).
We have already seen that $H_1$ is a solvable group of order 
$2p(p-1)$ generated by 
$\gamma_1$, $\gamma_2=\gamma_2(a)$, and $\gamma_3$, where 
$a\in F^\times$ is a primitive $2(p-1)-st$ root
of unity. As a consequence,
$|C_1|=|G\cdot P_1|=|G|/|H_1|=p+1$ 

Using $H_1=\langle \gamma_1,\gamma_2,\gamma_3\rangle $ and the 
explicit expressions for the $\gamma_i$, it can be checked
directly that no $g\in H_1$, $g\not= 1$, fixes any
$P\in C_2$. Therefore, $H_1\cap H_2=\{1\}$.
According to the proposition, 
the stabilizer $H_2$ of $P_2$ has order $p+1$.
This and $|G|=|H_1|\cdot |H_2|$ implies $G=H_1\cdot H_2$.
In other words, $H_2$ is a complement of $H_1$ in $G$.
(As {\it sets}, $G/H_1\cong PGL_2(p)/B\cong \ppp^1(F_0)$.)
In fact,
if $B$ denotes the (Borel) upper-triangular subgroup
of $PGL(2,p)$ then $H_1=\tau^{-1}(B)$.
Since $B$ is solvable and any abelian cover of a solvable group is
solvable, $H_1$ is solvable.
Since $B$ is not normal in $PGL_2(p)$, $H_1$ is not normal in $G$.

By the proposition, the divisor $D_2$ associated to $O_2$ 
has degree $2p(-1)>2g$, so by the Riemann-Roch theorem, 
${\rm dim}(L(D_2))=2p(p-1)+1-\frac{p-1}{2}$. 
The theorem implies that, in this case, the largest 
irreducible constituent of $L(D_2)$ is dimension
$d_G=p$.

\section{AG code constructions}


Throughout this section, we assume $X$, $G$, and $D$ are
as above and that $F$ is finite.

Let $P_1,...,P_n\in X(F)$ be distinct points
and $E=P_1+...+P_n\in Div(X)$ be stabilized by $G$. 
This implies that $G$ acts on the set
$supp(E)$ by permutation.
Assume ${\rm supp}(D)\cap {\rm supp}(E)=\emptyset$.
Let $C=C(D,E)$ denote the AG code
(\ref{eqn:AGcode}).
The group $G$ acts on $C$ by $g\in G$ sending
$c=(f(P_1),...,f(P_n))\in C$ to $c'=(f(g^{-1}(P_1)),...,f(g^{-1}(P_n)))$,
where $f\in L(D)$.
First, we observe that this map, denoted $\phi(g)$, is well-defined.
In other words, if $eval_E$ is not injective and
$c$ is also represented by $f'\in L(D)$, so
$c=(f'(P_1),...,f'(P_n))\in C$, then we can easily verify
$(f(g^{-1}(P_1)),...,f(g^{-1}(P_n)))=
(f'(g^{-1}(P_1)),...,f'(g^{-1}(P_n)))$.
(Indeed, $G$ acts on the set
$supp(E)$ by permutation.)
This map $\phi(g)$ induces a homomorphism of $G$ into the permutation
automorphism group of the code ${\rm Aut}(C)$, denoted 

\begin{equation}
\label{eqn:phi}
\phi:G\rightarrow {\rm Aut}(C) 
\end{equation}
(Prop. VII.3.3, \cite{Sti}).

The papers Wesemeyer \cite{W} and \cite{JT}, \cite{JK1} investigate 
$\phi$ in more detail.

\subsection{The kernel of $\phi$}

To investigate the kernel of this map $\phi$, we recall the following 
notion. Let $H\in Div(X)$ be any divisor.
We say that the space $L(H)$ {\bf separates points} if for
all points $P,Q\in X$,
$f(P)=f(Q)$ (for all $f\in L(H)$) 
implies $P=Q$ (see \cite{H}, chapter II, \S 7). 
Recall that if $D$ is ample then $L(D)$ separates points \cite{H}.

If $G$ is a group of automorphisms of $X$ defined over $F$
then $G$ induces an automorphism on the image of 
the evaluation map $eval_E : L(D)\rightarrow F^n$.
For this discussion, let us assume this is an injection.
(This is not a serious assumption.) 
To understand the kernel of this map $\phi$ in (\ref{eqn:phi}),
we'd like to know whether or not
$(f(P_1),...,f(P_n))=(f(g^{-1}P_1),...,f(g^{-1}P_n))$
implies $P_i=g^{-1}P_i$, for $1\leq i\leq n$.
In general, if $L(D)$ separates points then
\[
{\rm Ker}(\phi)=\{g\in G\ |\ g(P_i)=P_i,\ 1\leq i\leq n\}.
\]
It is known (proof of Prop. VII3.3, \cite{Sti}) that
if $n>2g+2$ then $\{g\in G\ |\ g(P_i)=P_i,\ 1\leq i\leq n\}$
is trivial. Therefore, if $n>2g+2$ and 
$D$ is ample then $\phi$ is injective.

\subsection{Permutation representations}

In this subsection, we show how theorems about AG codes
can, in some cases, give theorems about representations on
Riemann-Roch spaces.

Assume that $X/\fff$ is a hyperelliptic curve defined 
over a finite field $\fff$ 
of characteristic $p>2$ with automorphism 
group $G=Aut_\fff(X)$. 
Let $D$ be a $G$-equivariant
divisor on $X$, let ${\cal O}\subset X(\fff)$ be a 
$G$-orbit disjoint from the support of $D$, and let
$E=\sum_{P\in {\cal O}}P$. Let $P$ be the permutation automorphism
group of the code $C=C(D,E)$ defined in 
(\ref{eqn:AGcode}). 

Theorem 4.6 in \cite{W} implies that
if $n={\rm deg}(E)$ and
$t={\rm deg}(D)$ satisfy $n>{\rm max}(2t,2g+2)$ then
the map $\phi:G\rightarrow P$ is an isomorphism.
Using this, we regard $C$ as a $G$-module.
In particular, the (bijective) evaluation map 
$eval_E:L(D)\rightarrow C$ in (\ref{eqn:eval})
is $G$-equivariant.
Since $G$ acts (via its isomorphism with $P$) as a permutation
on $C$, we have proven the following result.

\begin{proposition}
Under the conditions above, the representation $\rho$ of $G$ on $L(D)$
is equivalent to a representation $\rho'$ with with property that,
for all $g\in G$, $\rho'(g)$ is a permutation matrix.

\end{proposition}

\subsection{Permutation decoding application}

If $C$ is an linear code with non-trivial permutation group
then this extra symmetry of the code may be useful in 
decoding.
Permutation decoding is discussed, for example, in 
Huffman and Pless \cite{HP}. It has also been implemented
in version 2.0 of the error-correcting computer
algebra package GUAVA \cite{GUAVA}. We recall briefly, for the
convenience of the reader, the main ideas.

We shall assume that $C$ is in standard form.
Let $C$ be a $[n,k,d]$ linear code over $GF(q)$,
let $t=[(d-1)/2]$, and let $G=(I_k,A)$ denote the generator matrix
in standard form. From this matrix $G$, it is well-known and easy to
show that one can compute
an encoder $E:GF(q)^k\rightarrow GF(q)^n$ with
image $C$, and a parity check matrix $H=(B,I_{n-k})$ 
in standard form, $B=-A^t$.

The key lemma is the following result: 
Suppose $v=c+e$, where $c\in C$ and $e\in GF(q)^n$
is an error vector with Hamming weight $wt(e)\leq t$. 
Under the above conditions,
the information symbols of $v$ are correct if and only if 
$wt(Hv)\leq t$. 

Let $P$ denote the permutation automorphism group of $C$. 
The permutation algorithm is:

\begin{enumerate}
\item
For each $p\in P$, compute $wt(H(pv))$ until one with
$wt(H(pv))\leq t$ is found (if none is found, the algorithm fails).

\item
Extract the information symbols from $pv$,
and use $E$ to compute codeword $c_p$ from them.

\item
Return $p^{-1}c_p={\rm Decode}(v)$. 
\end{enumerate}
For example, if $P$ acts transitively then 
permutation decoding will correct at least one error. 

The key problem is 
to find a set of permutations in $P$ which moves 
the non-zero positions in every possible error vector
of weight $\leq t$ out of the information positions.
(This set, called a {\it PD-set}, will be used in step
1 above instead of the entire set $P$.)

\begin{example}
We give two examples of MDS codes 
for which permutation decoding applies.

\begin{enumerate}
\item
This is an example of a $[7,3,5]$ one-point AG
code over $GF(7)$ arising from the hyperelliptic curve
$y^2=x^5-x$.

{\small{
\begin{verbatim}
p:=7;
F:=GF(p);
P<x>:=PolynomialRing(F);
f:=x^p-x;
X:=HyperellipticCurve(f);
Div := DivisorGroup(X);
Pls:=Places(X,1);
S:=[Pls[i] : i in [2..#Pls]];
m:=4;
D := m*(Div!Pls[1]);
AGC := AlgebraicGeometricCode(S, D);
Length(AGC);
Dimension(AGC);
MinimumDistance(AGC);
WeightDistribution(AGC);
PG := PermutationGroup(AGC);IdentifyGroup(PG);
ZP:=Center(PG);IdentifyGroup(PG/ZP);
IsTransitive(PG);
GeneratorMatrix(AGC);
\end{verbatim}
}}
This code has generator matrix in standard form given by

\[
G=
\left(
\begin{array}{ccccccc}
1 &0 &0 &2 &5 &1 &5\\
0 &1 &0 &1 &5 &5 &2\\
0 &0 &1 &5 &5 &2 &1
\end{array}
\right).
\]
Moreover, the permutation automorphism group of the code is a group
of order $42$ generated by

\[
S=\{ (1, 7)(2, 6)(3, 4), (1, 4, 5)(2, 6, 3),  (1, 3)(2, 4)(5, 6)\}.
\]
The elements of $S\, \cup \, S\cdot S$ can be used as a PD-set,
where $S\cdot S=\{s_1s_2\ |\ s_i\in S\}$.

\item
This is an example of a $[13,5,9]$ one-point AG
code over $GF(13)$ arising from the hyperelliptic curve
$y^2=x^{13}-x$. Similar MAGMA commands, but with $p=13$, 
yields that his code has generator matrix in standard form given by

\[
G=
\left(
\begin{array}{ccccccccccccc}
1 &  0 &  0 &  0 &  0 &  3 &  2 &  3 &  8 &  6 & 10 &  7 & 12\\
0 &  1 &  0 &  0 &  0 &  3 & 12 &  1 &  5 & 11 &  4 & 10  & 5\\
0 &  0 &  1 &  0 &  0 & 11 &  8  & 7 &  2  & 2 &  4 &  6 & 11\\
0 &  0 &  0 &  1 &  0 &  6 &  9 & 10 &  4 & 11 & 10 & 11  & 3\\
0 &  0 &  0 &  0 &  1 &  4 &  9 &  6 &  8 & 10 & 12 & 6 &  9
\end{array}
\right).
\]
Moreover, the permutation group is generated by

\[
\begin{array}{c}
\{p_1=(1, 2)(3, 8)(4, 12)(5, 7)(6, 9)(10, 11),\\
\ \ \ \ \ \ \ \ \ \ p_2=(2, 3, 4, 5, 6, 7, 8, 9, 10, 11, 12, 13)\}.
\end{array}
\]
We shall show that $P-\{1\}$ can be chosen to be
a PD-set for $\leq 4$ errors. The argument proceeds on a case-by-case basis.
One of the worst cases is when there are {\it errors 
in positions $1$, $2$, $12$ and $13$}. In this case, apply
(reading right-to-left) $p_1p_2p_1p_2^3$. This pushes the 
error from positions $(1,2,12,13)$ to $(8,11,13,6)$,
in particular out of the information positions\footnote{Additionally, 
this algorithm will, in {\it some}
cases, work even if there are $5$ errors (e.g., in positions $1,2,3,4,5$).
Because $d=9$, with $5$ errors we cannot be sure that the 
permutation decoded vector is the one which was sent.}.
\end{enumerate}

\end{example}

\begin{conjecture} For one-point AG codes $C$ associated
to $y^2=x^p-x$ over $GF(p)$ of length $n=p$,
permutation decoding always applies.
Its complexity is
at worst the size of the permutation group of $C$, 
which we conjecture to be $O(p^2)=O(g^2)=O(n^2)$.
\end{conjecture}

This matches the complexity of some known algorithms.

\begin{conjecture}
For one-point AG codes $C$ associated
to $y^2=x^p-x$ over $GF(p^2)$ of length $n=2p(p-1)+p$,
permutation decoding always applies and is more efficient
in terms of computational complexity than the standard
decoding algorithm in \cite{Sti}. We conjecture that, if the points
in $X(F)$ are arranged suitably then the image of the
$Aut_F(X)$ in the permutation group of $C$ may be used
as a PD-set. Its complexity is
at worst the size of the automorphism group of $X$, 
which is $O(p^2)=O(g^2)=O(n)$.
\end{conjecture}

If true, to our knowledge, this beats the complexity of other
decoding algorithms, such as those in \cite{Sti}.

\begin{example}
We give examples of two AG codes 
for which permutation decoding probably applies.

\begin{itemize}
\item
This is an example of a $[91,5,66]$ code constructed from 
the trace of a $[91,3,87]$ one-point AG
code over $F=GF(49)$ arising from the hyperelliptic curve
$y^2=x^7-x$. (We use the trace code only because MAGMA
version 2.10 cannot compute the permutation group of a
code over $GF(49)$.)  

{\small{
\begin{verbatim}
p:=7;
F:=GF(p^2);
P<x>:=PolynomialRing(F);
f:=x^p-x;
X:=HyperellipticCurve(f);
Div := DivisorGroup(X);
Pls:=Places(X,1);
S:=[Pls[i] : i in [2..#Pls]];
m:=4;
D := m*(Div!Pls[1]);
AGC := AlgebraicGeometricCode(S, D);
Length(AGC);
Dimension(AGC);
MinimumDistance(AGC);
WeightDistribution(AGC);
AGC0:=Trace(AGC,GF(p));
Length(AGC0);
Dimension(AGC0);
MinimumDistance(AGC0);
WeightDistribution(AGC0);
PG := PermutationGroup(AGC0);
#PG; 
ZP:=Center(PG);
#ZP; 
IsTransitive(PG);
GeneratorMatrix(AGC0);
\end{verbatim}
}}
The permutation group of the trace of the AG code is 
huge: $1073852196$ elements. The 
automorphism group $G=Aut_F(X)$, which has $672$ elements, 
of the curve acts on $X(F)$ with only two orbits,
$O_1=X(GF(7))$ of size $8$ and $O_2=X(F)-O_1$ of
size $84$. (This follows from Proposition \ref{prop:orbits}
but was verified using MAGMA in this case.)
Of course, in a practical application,
one would want to index the points of
$X(F)$ so that the information positions are contained in 
$O_2$.

\item
Let $E$ denote the sum of all the points in $O_2$
and let $D$ be the sum of all the points in $O_1$.
Note $E$ has degree $84$, $D$ degree $8$, and $X$ has genus $3$.
Therefore, by Theorem 4.6 in \cite{W},
the permutation automorphism group $P$ of the AG code $C=C(D,E)$  
satisfies $P\cong G$. In other words, the map
$\phi$ in (\ref{eqn:phi}) is injective (which also
follows from the discussion above) and surjective.
\end{itemize}
\end{example}


\begin{thebibliography}{99}  

%\bibitem[Bo1]{Bo1} N. Borne, {\it Une formule de Riemann-Roch
%equivariante pour des courbes,} preprint
%1999
%
%\bibitem[Bo2]{Bo2} -----, {\it Structure de groupe de Grothendieck
%\'equivariant d'une courbe et modules galoisiens,} preprint
%
%\bibitem[Bo3]{Bo3} -----,{\it Modules galoisiens sur les courbes: 
%une introduction,} S\'em. et Congr\`es, SMF \underline{5}(2001)147-159.
%
%\bibitem[B]{B} T. Breuer, {Characters and automorphism
%groups of Riemann surfaces}, London Math. Soc.
%Lecture Notes, 1999.
%
%\bibitem[CR]{CR} C. Curtis, I. Reiner, {Methods of 
%representation theory, I}, Wiley-Interscience, 1981.
% 
%\bibitem[F]{F} W. Fulton, {Algebraic curves},
%Benjamin, 1969.
%
%\bibitem[FH]{FH} ------ and J. Harris, {Representation theory:
%a first course}, Springer-Verlag, 1991.

%\bibitem[GAP]{GAP4}
%The GAP~Group, \emph{GAP -- Groups, Algorithms, and Programming, 
%Version 4.3}; 2002, (\verb+http://www.gap-system.org+).

\bibitem[G]{G} N. G\"ob, ``Computing the Automorphism Groups of
Hyperelliptic Function Fields,'' available at
\verb+http://front.math.ucdavis.edu/math.NT/0305284+.

\bibitem[GUAVA]{GUAVA}
GUAVA package for GAP 4.4 (\verb+http://www.gap-system.org+),
available at
\verb+http://cadigweb.ew.usna.edu/~wdj/gap/GUAVA/+
or
\newline
\verb+http://www.gap-system.org/Packages/guava.html+

\bibitem[H]{H} R. Hartshorne, {Algebraic geometry},
Springer-Verlag, 1977.

\bibitem[HP]{HP} W. C. Huffman and V. Pless, {Fundamentals of 
error-correcting codes}, Cambridge Univ. Press, 2003.

\bibitem[JK1]{JK1} D. Joyner and A. Ksir, 
``Automorphism groups of some AG codes'' available at
\verb+http://cadigweb.ew.usna.edu/~wdj/papers.html+

\bibitem[JK2]{JK2} D. Joyner and A. Ksir, 
``Representations of finite groups on Riemann-Roch
spaces, II'' available at
\verb+http://front.math.ucdavis.edu/math.AG/0312383+

\bibitem[JT]{JT} D. Joyner and W. Traves, 
``Representations of finite groups on Riemann-Roch
spaces'' available at
\verb+http://front.math.ucdavis.edu/math.AG/0210408+

%\bibitem[Ka]{Ka} E. Kani, {\it The Galois-module structure of the 
%space of holomorphic differential forms on a curve,}
%J. reine angew. Math \underline{367}(1986)187-206.
%
%\bibitem[K]{K} B. K\"ock, {\it Galois structure of Zariski cohomology 
%for weakly ramified covers of curves,} 
%math.AG/0207124, available at
%\verb+http://front.math.ucdavis.edu/+
%
%\bibitem[L]{L} D. Lorenzini, {An invitation to arithmetic 
%geometry}, Grad. Studies in Math, AMS, 1996.

\bibitem[MAGMA]{M} 
W. Bosma, J. Cannon, C. Playoust,
{\it The MAGMA algebra system, I: The user language},
J. Symb. Comp., \underline{24}(1997)235-265.
\newline
(\verb+http://www.maths.usyd.edu.au:8000/u/magma/+).

%\bibitem[MP]{MP} I. Morrison and H. Pinkham,
%{\it Galois Weierstass points and Hurwitz characters},
%Annals of Math., \underline{124}(1986)591-625.
%
%\bibitem[N]{N} S. Nakajima, {\it Galois module structure of
%cohomology groups of an algebraic variety},
%Inv. Math. \underline{75}(1984)1-8
%
%\bibitem[Se]{Se} J.-P. Serre, {Linear representations of 
%finite groups}, Springer-Verlag, 1977.
%
%\bibitem[St]{St} S. Stepanov, {Codes on algebraic curves},
%Kluwer, NY, 1999.

\bibitem[Sti]{Sti} H. Stichtenoth, {Algebraic function
fields and codes}, Springer-Verlag, 1993.

\bibitem[W]{W} S. Wesemeyer, {\it On the Automorphism Group of 
Various Goppa Codes},
IEEE Trans. Info. Theory., 44(1998)630-643.
\end{thebibliography}

\end{document}
